\section{Proofs of the normalization-risk comparisons}
\label{sec:comparison-proofs}

This appendix gives the finite-moment formulas, approximation bounds, and
interval-certificate construction for Theorems~\ref{thm:conference-c1}
and~\ref{thm:conference-c2}. All calculations use the binary model,
fresh sampling, and population preprocessing of
Section~\ref{sec:conference-setting}. The implementation is
\texttt{experiments/check\_normalization\_risk\_theorem.py}; inside the
locked project environment, \texttt{make check-normalization-risk} runs
both the certificates and independent numerical checks. Packaging an
anonymized reproducibility archive remains a submission task.

\subsection{Exact moments and the variance decomposition}

Write $D=\eta P=\diag(d_C,d_R)$ and
$s_k=(V_{CC,k},V_{CR,k},V_{RR,k})^\top$. Set
\[
 t_B=c^2+\frac{ab-c^2}{B},\qquad
 m_B=c\left(b+\frac{1-4b}{B}\right).
\]
The exact recurrence is $s_{k+1}=L_Ds_k+\sigma^2g_D$, where
\[
 L_D=\begin{pmatrix}
 (1-d_Ca)^2 & -2d_C(1-d_Ca)c & d_C^2t_B\\
 -d_R(1-d_Ca)c & (1-d_Ca)(1-d_Rb)+d_Cd_Rt_B & -d_Cc+d_Cd_Rm_B\\
 d_R^2t_B & -2d_Rc+2d_R^2m_B & (1-d_Rb)^2+d_R^2b(1-4b)/B
 \end{pmatrix},
\]
\[
 g_D=\frac1B(d_C^2a,d_Cd_Rc,d_R^2b)^\top,\qquad
 s_0=((\beta_C^*)^2,\beta_C^*\beta_R^*,(\beta_R^*)^2)^\top.
\]
To verify this, let $\widehat C$ be the averaged batch second moment.
The common squared coordinate is the constant $a$, and the binary law gives
\[
 \E x_R^4=b(1-3b),\quad \E x_C^2x_R^2=ab,\quad
 \E x_Cx_R^3=c(1-3b).
\]
Consequently, $\E\widehat c^2=t_B$,
$\E(\widehat c\widehat b)=m_B$, and
$\E\widehat b^2=b^2+b(1-4b)/B$.
Expanding $\E[(I-D\widehat C)V_k(I-\widehat CD)]$ yields $L_D$.
The conditional mean-zero noise assumption removes cross terms, and the
noise injection has covariance $\sigma^2DCD/B$. Fresh sampling makes the
new update independent of the previous error.

For \eqref{eq:conference-decomposition}, run the noiseless recursion on
the same feature stream. Its mean error is
$e_k=-(I-DC)^k\beta^*$ and its covariance is $A_k\succeq0$.
Conditional on all sampled features, the noise-driven component has mean
zero, so its cross moment with the noiseless component vanishes.
Its covariance is $\sigma^2N_k\succeq0$. This proves the decomposition.
Equivalently, if $\mathcal T_D(V)=\E[(I-D\widehat C)V(I-\widehat CD)]$,
then
\[
 V_k=\mathcal T_D^k(\beta^*\beta^{*\top})+
 \frac{\sigma^2}{B}\sum_{j=0}^{k-1}\mathcal T_D^j(DCD).
\]

\subsection{A scalar approximation with a controlled remainder}

Under independence, the diagonal moments close under
\[
 K_D=\begin{pmatrix}A&e\\f&D_0\end{pmatrix},\quad
 A=(1-h)^2,\quad e=\frac{d_C^2ab}{B},\quad f=\frac{d_R^2ab}{B},\quad
 D_0=(1-d_Rb)^2+\frac{d_R^2b(1-4b)}B.
\]
Here $D_0$ is a scalar, distinct from $D$. All entries are nonnegative.
Define
\[
 w=\tfrac12(a,(1-q)^2)^\top,\quad
 v_0=((\beta_C^*)^2,(\beta_R^*)^2)^\top,\quad
 g=(d_C^2a,d_R^2b)^\top/B,
\]
and let $E_0=\max(e,f)$ and $\|K_D\|_\infty\leq\rho<1$.
For $\phi_k(z)=\sum_{j=0}^{k-1}z^j$, define
\[
 \widetilde R_{R,k}=
 w_C\{A^k(\beta_C^*)^2+\sigma^2g_C\phi_k(A)\}
 +w_R\{D_0^k(\beta_R^*)^2+\sigma^2g_R\phi_k(D_0)\}.
\]
For every $k\geq1$,
\begin{equation}
 0\leq R_{R,k}-\widetilde R_{R,k}\leq E_F+\sigma^2E_N,\qquad
 E_F=\|w\|_1\|v_0\|_\infty E_0k\rho^{k-1},\quad
 E_N=\frac{\|w\|_1\|g\|_\infty E_0}{(1-\rho)^2}.
 \label{eq:comparison-remainder}
\end{equation}
Indeed, with $K_0=\diag(A,D_0)$ and $E=K_D-K_0\geq0$ entrywise,
\[
 K_D^k-K_0^k=\sum_{j=0}^{k-1}K_D^{k-1-j}EK_0^j\geq0.
\]
Its infinity norm is at most $E_0k\rho^{k-1}$. Multiplication by $w$
and $v_0$ proves the homogeneous bound. Applying the same identity to the
noise sum and using $\sum_{j=1}^{k-1}j\rho^{j-1}\leq(1-\rho)^{-2}$
proves the noise bound.

The common terms cancel between the U and S proxies. Write
\[
 \widetilde F=w_R(\beta_R^*)^2(D_S^k-D_U^k),\qquad
 \widetilde N=w_R\{g_{R,S}\phi_k(D_S)-g_{R,U}\phi_k(D_U)\},
\]
where $D_S,D_U$ are the respective scalar $D_0$ values. For the exact
contrast $F+\sigma^2N$, subtracting the two nonnegative remainders gives
\[
 F_-:=\widetilde F-E_{F,U}\leq F\leq\widetilde F+E_{F,S}=:F_+,
\]
\[
 N_-:=\widetilde N-E_{N,U}\leq N\leq\widetilde N+E_{N,S}=:N_+.
\]
When $F_+<0<N_-$, division by positive denominators gives
\eqref{eq:conference-threshold}. In particular,
$F_++\sigma^2N_+<0$ guarantees lower rare risk for S, whereas
$F_-+\sigma^2N_->0$ guarantees lower rare risk for U. The homogeneous
proxy already includes rare-feature update variance through $D_0$;
it is not simply the squared mean error.

\subsection{Certificate for Theorem~\ref{thm:conference-c1}}

Partition each of the $q$ and $h$ intervals into eight equal closed
pieces. On each of the resulting 64 rectangles, evaluate the preceding
bounds with interval arithmetic and
$(\beta_C^*)^2,(\beta_R^*)^2\in[.9025,1.1025]$.
For each configuration use upper bounds on
$\rho=\max(A+e,D_0+f)$ and $E_0=\max(e,f)$.
The certificate obtains $\rho<.999719$ throughout. The remaining
cross moment propagates by
\[
 \gamma=(1-d_Ca)(1-d_Rb)+d_Cd_Rab/B
\]
with zero noise injection. The certificate checks $0\leq\gamma<1$
on every rectangle. Thus the full homogeneous second-moment map is
stable, and the fixed finite noise injection gives bounded noisy moments.

Every rectangle satisfies $F_+<0<N_-$. The hulls of the resulting
$F+\sigma^2N$ intervals give Table~\ref{tab:conference-c1}. The respective
noise-threshold enclosures, in its row order, are
$[10.46,17.76]$, $[.53,1.405]$, $[112.61,184.22]$, and $[5.69,14.49]$.
They imply the four signs uniformly whenever
$1.405<\sigma^2<5.69$. The narrower interval in the theorem gives its
displayed risk-gap bounds. Closed rectangles cover their shared boundaries.

\subsection{Outward-rounded arithmetic}

Each interval endpoint is represented by an integer multiple of
$10^{-60}$. Inputs are rational strings or integers. Addition is exact;
products use the minimum and maximum of all four endpoint products and
round outward using integer division. If $[l,u]$ stores the integer
endpoints of a positive denominator interval, its reciprocal is enclosed by
\[
 \left[\frac{\lfloor10^{120}/u\rfloor}{10^{60}},
 \frac{\lceil10^{120}/l\rceil}{10^{60}}\right].
\]
Negative denominators are reduced to this case. Powers and finite sums
use repeated squaring and addition. Induction over the arithmetic
operations gives enclosure of each exact rational expression. All
acceptance checks compare integer endpoints, including containment in
the reported decimal bounds. Floating-point conversions are used only
for display. Thus the parameter boxes are covered by finite enclosure
calculations, not by testing isolated parameter values.

\subsection{Proof of Theorem~\ref{thm:conference-c2}}

For stability, a feature vector with $\|z\|^2\leq L$ satisfies
\[
 \E[(I-\eta zz^\top)^2]\preceq I-\eta(2-\eta L)H,
 \qquad H=\E[zz^\top].
\]
If $H\succ0$ and $0<\eta<2/L$, this contracts expected squared
error. For U, take $L=.25+.98^2=1.2104$, $\eta=.04$, and use
$\lambda_{\min}(C)\geq.0196-.01=.0096$. For S, $L=50$,
$\eta=.01$, and
\[
 H=\begin{pmatrix}1&\ell/7\\\ell/7&1\end{pmatrix},\qquad
 \lambda_{\min}(H)\geq6/7.
\]
These sufficient stability bounds hold throughout the claimed interval.

For the risks, augment $L_D$ by the unit-noise column $g_D$ and a final
row $(0,0,0,1)$. Its $100$th power gives the homogeneous map and finite
noise sum. The aggregate and rare risk weights on $s_k$ are
\[
 w_A=(a/2,c,b/2)^\top,\qquad
 w_R=(a/2,\ell(1-q)/2,(1-q)^2/2)^\top.
\]
Partition $[.9,1]$ into 128 equal closed intervals. On each interval,
the same integer arithmetic encloses the two augmented powers, their
action on $s_0=(4,2,1)^\top$, and the contrasts
$F_W+\sigma^2N_W$ for $\sigma^2\in[3.8,4.1]$.
Signed matrix entries use all four endpoint products; no entrywise
positivity is assumed. The union satisfies both risk enclosures in the
theorem. Before risk weighting, the same powers certify positive
diagonal contrasts and a negative cross-moment contrast throughout.
This proves the coefficient MSE inequalities as well.

\subsection{Checks distinct from the certificates}

Independent ordered-batch and short-path enumeration checks the
three-moment formulas against actual gradient updates. Comparisons
with a separate full moment operator cover 64 endpoints, including
mixed-sign signals, dependence boundaries, an interior law, and
$B=1,2$. These floating-point checks validate implementations; the
interval arithmetic above supplies the continuous-box certificates.
Neither form of verification establishes priority or robustness to
procedures outside the stated comparison.
